\documentclass[12pt]{article}
\usepackage[a4paper,top=2.5cm,bottom=2.5cm,left=2cm,right=2cm]{geometry}

% Packages de base
\usepackage[utf8]{inputenc}
\usepackage[T1]{fontenc}
\usepackage[french]{babel} % date et titres en français
\usepackage{lipsum}        % texte d'exemple
\usepackage{fancyhdr}      % pour personnaliser en-têtes et pieds

% Réglages fancyhdr
\pagestyle{fancy}
\fancyhf{} % vide tout
\newcommand{\titrechapitre}{Algèbre linéaire}
\newcommand{\numerochapitre}{X}

\renewcommand{\headrulewidth}{0.4pt}
\renewcommand{\footrulewidth}{0.4pt}
\fancyfoot[L]{Notes de Raphaël JONTEF} 
\fancyfoot[C]{\thepage} 
\fancyfoot[R]{Enseirb-Matmeca}
\fancyhead[L]{Cours de Mathématiques}
\fancyhead[R]{\titrechapitre}


\begin{document}
% Page de titre custom
\begin{titlepage}
  \centering
  \vspace*{3cm}
  {\Huge\bfseries Chapitre \numerochapitre \\[1ex] \titrechapitre\par}
  \vspace{2cm}
  {\Large Notes rédigées par Raphaël JONTEF\par}
  {\Large Cours enseigné par François DUFOUR\par}
  \vspace{4cm}
  {\Big Enseirb-Matmeca\par}
  {\Big filière informatique\par}

  \vfill
  {\large \today\par}
\end{titlepage}

\tableofcontents % plan du cours
\newpage

% -----------Corps du document--------------

\section{Introduction}
\lipsum[1]

\section{Deuxième partie}
\lipsum[2]

\section{Troisième partie}
\lipsum[3] \lipsum[4]

\end{document}
